\documentclass[11pt]{article}

\usepackage[margin=1in]{geometry}
\usepackage[T1]{fontenc}
\usepackage[utf8]{inputenc}
\usepackage{lmodern}
\usepackage{amsmath,amssymb}
\usepackage{booktabs}
\usepackage{enumitem}
\usepackage{hyperref}
\usepackage{tabularx}
\usepackage{xcolor}

\hypersetup{
  colorlinks=true,
  linkcolor=blue!50!black,
  urlcolor=blue!50!black,
  citecolor=blue!50!black,
  pdftitle={Promising Directions Assessment for p-vs-np}
}

\title{Promising Directions Assessment for \texttt{p-vs-np}}
\author{Neuro-symbolic AI Group\\Idiap Research Institute}
\date{March 16, 2026}

\begin{document}
\maketitle

\begin{abstract}
This report identifies the most promising mathematical directions currently surfaced by the framework for
\texttt{p-vs-np}. The assessment is content-first: each route is described by
its theorem pair, typed bridge object, preserved invariant, missing route lemma, and failure logic. The novelty claim
is assessed at the level of mathematical program rather than theorem announcement. The conclusion is that the current
run contains real route-level novelty and a plausible discovery direction, but not a validated theorem-level discovery.
\end{abstract}

\section{Run Context}
\begin{itemize}[leftmargin=*,itemsep=0.25em]
\item Generated (UTC): 2026-03-16T22:26:32.766903+00:00
\item Source run: 2026-03-16T22:26:28.893056+00:00
\item Evidence count: 0
\item Unresolved hypotheses: 22
\end{itemize}

\section{Executive Summary}
There is real programmatic novelty in the sense of a new content-level bridge program: the framework isolates specific theorem-pair deltas, proposes typed intermediate carriers, and exposes missing route lemmas. There is not yet evidence of theorem-level novelty, because the decisive lemmas remain open and the run has not accumulated external or adversarial evidence.

\section{Most Promising Directions}
\subsection{1. MCSP witness family: major consequences such as NP not in P/poly}
\begin{itemize}[leftmargin=*,itemsep=0.25em]
\item Bridge ID: \texttt{bridge.001}
\item Type: magnification bridge
\item Priority: 0.744
\item Route: \texttt{st.def.hardness-magnification-phenomenon-programme-leve} $\rightarrow$ \texttt{st.thm.oliveira-pich-santhanam-magnification-near-state}
\item Abstraction: MCSP witness family
\item Carrier: family of instances equipped with a monotone witness measure
\item Invariant: MCSP complexity profile in nonuniform-circuits
\item Route synopsis: Move from 'st.def.hardness-magnification-phenomenon-programme-leve' on problem 'MCSP' in model 'nonuniform-circuits' using 'A barely-nontrivial lower bound for an explicit problem Q in a weak model' toward 'st.thm.oliveira-pich-santhanam-magnification-near-state' in model 'nonuniform-circuits' to obtain 'major consequences such as NP not in P/poly'.
\end{itemize}
\paragraph{Why It Is Promising}
\begin{itemize}[leftmargin=*,itemsep=0.2em]
\item Bridge priority 0.744 with sketch score 0.905, mathematical specificity 1.000, evidence signal 0.188, and sketch operationalization 1.000.
\item The route synopsis is: Move from 'st.def.hardness-magnification-phenomenon-programme-leve' on problem 'MCSP' in model 'nonuniform-circuits' using 'A barely-nontrivial lower bound for an explicit problem Q in a weak model' toward 'st.thm.oliveira-pich-santhanam-magnification-near-state' in model 'nonuniform-circuits' to obtain 'major consequences such as NP not in P/poly'.
\item The proof family 'invariant-lift-through-intermediate-object' is used because Introduce a conjectured intermediate object or invariant, prove that the source regime induces it, then force the target consequence through that invariant.
\item Linked hypotheses: 4 with verdict profile \{'supported': 3, 'inconclusive': 1\} and best confidence 0.804.
\end{itemize}
\paragraph{Content-wise Novelty Argument}
\begin{itemize}[leftmargin=*,itemsep=0.2em]
\item The route is theorem-pair specific: it tries to move from 'A barely-nontrivial lower bound for an explicit problem Q in a weak model' to 'major consequences such as NP not in P/poly' rather than proposing a generic topic-level bridge.
\item The candidate novelty is the object 'MCSP witness family', typed as 'family of instances equipped with a monotone witness measure', which is meant to make the hidden transport mechanism mathematically explicit.
\item The route is constrained by the preserved invariant 'MCSP complexity profile in nonuniform-circuits', so the proposal is not only to reach the target consequence but to preserve a specific mathematical profile while doing so.
\item The novelty remains programmatic rather than theorem-level because the decisive missing lemma is still: Prove that a hardness-magnification witness family turns 'A barely-nontrivial lower bound for an explicit problem Q in a weak model' into 'major consequences such as NP not in P/poly' while preserving 'MCSP complexity profile in nonuniform-circuits'.
\item The route includes an explicit falsification probe: Allow 'MCSP witness family' to ignore the source-side content 'A barely-nontrivial lower bound for an explicit problem Q in a weak model' and only preserve topical overlap. bridge-degrades-or-fails The bridge object must absorb the actual source-side mathematical content rather than merely name the antecedent.
\end{itemize}
\paragraph{Key Obligations}
\begin{itemize}[leftmargin=*,itemsep=0.2em]
\item Prove the source-to-object claim: the source-side content 'A barely-nontrivial lower bound for an explicit problem Q in a weak model' available at 'st.def.hardness-magnification-phenomenon-programme-leve' can be organized into 'MCSP witness family' as a family of instances equipped with a monotone witness measure.
\item Prove the payoff claim: once 'MCSP witness family' carries the route data from 'A barely-nontrivial lower bound for an explicit problem Q in a weak model' toward 'major consequences such as NP not in P/poly', it forces 'NP-not-in-P/poly' at 'st.thm.oliveira-pich-santhanam-magnification-near-state'.
\item Prove the transport claim: every admissible realization of 'MCSP witness family' preserves 'MCSP complexity profile in nonuniform-circuits' while crossing the route nonuniform-circuits.
\item Discharge the route lemma: Prove that a hardness-magnification witness family turns 'A barely-nontrivial lower bound for an explicit problem Q in a weak model' into 'major consequences such as NP not in P/poly' while preserving 'MCSP complexity profile in nonuniform-circuits'.
\end{itemize}
\paragraph{Main Risks}
\begin{itemize}[leftmargin=*,itemsep=0.2em]
\item Failure mode: 'MCSP witness family' never becomes more than topical overlap because the route-specific lemma remains missing. Prove that a hardness-magnification witness family turns 'A barely-nontrivial lower bound for an explicit problem Q in a weak model' into 'major consequences such as NP not in P/poly' while preserving 'MCSP complexity profile in nonuniform-circuits'.
\item Abort logic: Abort if the negative-route intervention does not expose the expected failure: Allow 'MCSP witness family' to ignore the source-side content 'A barely-nontrivial lower bound for an explicit problem Q in a weak model' and only preserve topical overlap. => bridge-degrades-or-fails.
\item Global limitation: evidence\_count=0 and unresolved hypotheses=22, so the route is still under mathematical pressure.
\end{itemize}
\paragraph{Expert Actions}
\begin{itemize}[leftmargin=*,itemsep=0.2em]
\item Discharge the route-specific lemma: Prove that a hardness-magnification witness family turns 'A barely-nontrivial lower bound for an explicit problem Q in a weak model' into 'major consequences such as NP not in P/poly' while preserving 'MCSP complexity profile in nonuniform-circuits'.
\item Stabilize the main bridge lemma: Source-to-object lemma: the source-side content 'A barely-nontrivial lower bound for an explicit problem Q in a weak model' at 'st.def.hardness-magnification-phenomenon-programme-leve' yields a realization of 'MCSP witness family' as a family of instances equipped with a monotone witness measure, satisfying admits source-side witness data for 'A barely-nontrivial lower bound for an explicit problem Q in a weak model'; preserves the route invariant 'MCSP complexity profile in nonuniform-circuits'; forces the target-side payoff 'NP-not-in-P/poly'.
\item Run the designated counterfactual probe: Allow 'MCSP witness family' to ignore the source-side content 'A barely-nontrivial lower bound for an explicit problem Q in a weak model' and only preserve topical overlap. bridge-degrades-or-fails The bridge object must absorb the actual source-side mathematical content rather than merely name the antecedent.
\end{itemize}
\subsection{2. MCSP witness family: major lower bounds for NP in much stronger models (including superpolynomial formula/circuit lower bounds, and in some parameter...}
\begin{itemize}[leftmargin=*,itemsep=0.25em]
\item Bridge ID: \texttt{bridge.002}
\item Type: magnification bridge
\item Priority: 0.732
\item Route: \texttt{st.def.hardness-magnification-phenomenon-programme-leve} $\rightarrow$ \texttt{st.thm.oliveira-santhanam-hardness-magnification-for-na}
\item Abstraction: MCSP witness family
\item Carrier: family of instances equipped with a monotone witness measure
\item Invariant: MCSP complexity profile in nonuniform-circuits
\item Route synopsis: Move from 'st.def.hardness-magnification-phenomenon-programme-leve' on problem 'MCSP' in model 'nonuniform-circuits' using 'A barely-nontrivial lower bound for an explicit problem Q in a weak model' toward 'st.thm.oliveira-santhanam-hardness-magnification-for-na' in model 'formula-size' to obtain 'major lower bounds for NP in much stronger models (including superpolynomial formula/circuit lower bounds, and in some parameter...'.
\end{itemize}
\paragraph{Why It Is Promising}
\begin{itemize}[leftmargin=*,itemsep=0.2em]
\item Bridge priority 0.732 with sketch score 0.897, mathematical specificity 1.000, evidence signal 0.188, and sketch operationalization 1.000.
\item The route synopsis is: Move from 'st.def.hardness-magnification-phenomenon-programme-leve' on problem 'MCSP' in model 'nonuniform-circuits' using 'A barely-nontrivial lower bound for an explicit problem Q in a weak model' toward 'st.thm.oliveira-santhanam-hardness-magnification-for-na' in model 'formula-size' to obtain 'major lower bounds for NP in much stronger models (including superpolynomial formula/circuit lower bounds, and in some parameter...'.
\item The proof family 'invariant-lift-through-intermediate-object' is used because Introduce a conjectured intermediate object or invariant, prove that the source regime induces it, then force the target consequence through that invariant.
\item Linked hypotheses: 4 with verdict profile \{'supported': 3, 'inconclusive': 1\} and best confidence 0.804.
\end{itemize}
\paragraph{Content-wise Novelty Argument}
\begin{itemize}[leftmargin=*,itemsep=0.2em]
\item The route is theorem-pair specific: it tries to move from 'A barely-nontrivial lower bound for an explicit problem Q in a weak model' to 'major lower bounds for NP in much stronger models (including superpolynomial formula/circuit lower bounds, and in some parameter...' rather than proposing a generic topic-level bridge.
\item The candidate novelty is the object 'MCSP witness family', typed as 'family of instances equipped with a monotone witness measure', which is meant to make the hidden transport mechanism mathematically explicit.
\item The route is constrained by the preserved invariant 'MCSP complexity profile in nonuniform-circuits', so the proposal is not only to reach the target consequence but to preserve a specific mathematical profile while doing so.
\item The novelty remains programmatic rather than theorem-level because the decisive missing lemma is still: Prove that a hardness-magnification witness family preserves 'MCSP complexity profile in nonuniform-circuits' while transporting the route from the nonuniform-circuits regime to the formula-size regime strongly enough to derive 'major lower bounds for NP in much stronger models (including superpolynomial formula/circuit lower bounds, and in some parameter...'.
\item The route includes an explicit falsification probe: Allow 'MCSP witness family' to ignore the source-side content 'A barely-nontrivial lower bound for an explicit problem Q in a weak model' and only preserve topical overlap. bridge-degrades-or-fails The bridge object must absorb the actual source-side mathematical content rather than merely name the antecedent.
\end{itemize}
\paragraph{Key Obligations}
\begin{itemize}[leftmargin=*,itemsep=0.2em]
\item Prove the source-to-object claim: the source-side content 'A barely-nontrivial lower bound for an explicit problem Q in a weak model' available at 'st.def.hardness-magnification-phenomenon-programme-leve' can be organized into 'MCSP witness family' as a family of instances equipped with a monotone witness measure.
\item Prove the payoff claim: once 'MCSP witness family' carries the route data from 'A barely-nontrivial lower bound for an explicit problem Q in a weak model' toward 'major lower bounds for NP in much stronger models (including superpolynomial formula/circuit lower bounds, and in some parameter...', it forces 'NP-not-in-P/poly' at 'st.thm.oliveira-santhanam-hardness-magnification-for-na'.
\item Prove the transport claim: every admissible realization of 'MCSP witness family' preserves 'MCSP complexity profile in nonuniform-circuits' while crossing the route nonuniform-circuits $\rightarrow$  formula-size.
\item Discharge the route lemma: Prove that a hardness-magnification witness family preserves 'MCSP complexity profile in nonuniform-circuits' while transporting the route from the nonuniform-circuits regime to the formula-size regime strongly enough to derive 'major lower bounds for NP in much stronger models (including superpolynomial formula/circuit lower bounds, and in some parameter...'.
\end{itemize}
\paragraph{Main Risks}
\begin{itemize}[leftmargin=*,itemsep=0.2em]
\item Failure mode: 'MCSP witness family' never becomes more than topical overlap because the route-specific lemma remains missing. Prove that a hardness-magnification witness family preserves 'MCSP complexity profile in nonuniform-circuits' while transporting the route from the nonuniform-circuits regime to the formula-size regime strongly enough to derive 'major lower bounds for NP in much stronger models (including superpolynomial formula/circuit lower bounds, and in some parameter...'.
\item Abort logic: Abort if the negative-route intervention does not expose the expected failure: Allow 'MCSP witness family' to ignore the source-side content 'A barely-nontrivial lower bound for an explicit problem Q in a weak model' and only preserve topical overlap. => bridge-degrades-or-fails.
\item Global limitation: evidence\_count=0 and unresolved hypotheses=22, so the route is still under mathematical pressure.
\end{itemize}
\paragraph{Expert Actions}
\begin{itemize}[leftmargin=*,itemsep=0.2em]
\item Discharge the route-specific lemma: Prove that a hardness-magnification witness family preserves 'MCSP complexity profile in nonuniform-circuits' while transporting the route from the nonuniform-circuits regime to the formula-size regime strongly enough to derive 'major lower bounds for NP in much stronger models (including superpolynomial formula/circuit lower bounds, and in some parameter...'.
\item Stabilize the main bridge lemma: Source-to-object lemma: the source-side content 'A barely-nontrivial lower bound for an explicit problem Q in a weak model' at 'st.def.hardness-magnification-phenomenon-programme-leve' yields a realization of 'MCSP witness family' as a family of instances equipped with a monotone witness measure, satisfying admits source-side witness data for 'A barely-nontrivial lower bound for an explicit problem Q in a weak model'; preserves the route invariant 'MCSP complexity profile in nonuniform-circuits'; forces the target-side payoff 'NP-not-in-P/poly'.
\item Run the designated counterfactual probe: Allow 'MCSP witness family' to ignore the source-side content 'A barely-nontrivial lower bound for an explicit problem Q in a weak model' and only preserve topical overlap. bridge-degrades-or-fails The bridge object must absorb the actual source-side mathematical content rather than merely name the antecedent.
\end{itemize}
\subsection{3. MCSP witness family: There exist general-purpose magnification theorems (not problem-specific) showing that certain mild lower bounds for...}
\begin{itemize}[leftmargin=*,itemsep=0.25em]
\item Bridge ID: \texttt{bridge.016}
\item Type: magnification bridge
\item Priority: 0.711
\item Route: \texttt{st.def.hardness-magnification-phenomenon-programme-leve} $\rightarrow$ \texttt{st.thm.atserias-m-uller-simple-general-magnification-fr}
\item Abstraction: MCSP witness family
\item Carrier: family of instances equipped with a monotone witness measure
\item Invariant: exponential-lower-bound witness complexity
\item Route synopsis: Move from 'st.def.hardness-magnification-phenomenon-programme-leve' on problem 'MCSP' in model 'nonuniform-circuits' using 'A barely-nontrivial lower bound for an explicit problem Q in a weak model' toward 'st.thm.atserias-m-uller-simple-general-magnification-fr' to obtain 'There exist general-purpose magnification theorems (not problem-specific) showing that certain mild lower bounds for...'.
\end{itemize}
\paragraph{Why It Is Promising}
\begin{itemize}[leftmargin=*,itemsep=0.2em]
\item Bridge priority 0.711 with sketch score 0.890, mathematical specificity 1.000, evidence signal 0.200, and sketch operationalization 1.000.
\item The route synopsis is: Move from 'st.def.hardness-magnification-phenomenon-programme-leve' on problem 'MCSP' in model 'nonuniform-circuits' using 'A barely-nontrivial lower bound for an explicit problem Q in a weak model' toward 'st.thm.atserias-m-uller-simple-general-magnification-fr' to obtain 'There exist general-purpose magnification theorems (not problem-specific) showing that certain mild lower bounds for...'.
\item The proof family 'invariant-lift-through-intermediate-object' is used because Introduce a conjectured intermediate object or invariant, prove that the source regime induces it, then force the target consequence through that invariant.
\item Linked hypotheses: 5 with verdict profile \{'supported': 4, 'inconclusive': 1\} and best confidence 0.804.
\end{itemize}
\paragraph{Content-wise Novelty Argument}
\begin{itemize}[leftmargin=*,itemsep=0.2em]
\item The route is theorem-pair specific: it tries to move from 'A barely-nontrivial lower bound for an explicit problem Q in a weak model' to 'There exist general-purpose magnification theorems (not problem-specific) showing that certain mild lower bounds for...' rather than proposing a generic topic-level bridge.
\item The candidate novelty is the object 'MCSP witness family', typed as 'family of instances equipped with a monotone witness measure', which is meant to make the hidden transport mechanism mathematically explicit.
\item The route is constrained by the preserved invariant 'exponential-lower-bound witness complexity', so the proposal is not only to reach the target consequence but to preserve a specific mathematical profile while doing so.
\item The novelty remains programmatic rather than theorem-level because the decisive missing lemma is still: Prove that a hardness-magnification witness family upgrades the route from major separation consequence to exponential lower bound without losing 'exponential-lower-bound witness complexity'.
\item The route includes an explicit falsification probe: Allow 'MCSP witness family' to ignore the source-side content 'A barely-nontrivial lower bound for an explicit problem Q in a weak model' and only preserve topical overlap. bridge-degrades-or-fails The bridge object must absorb the actual source-side mathematical content rather than merely name the antecedent.
\end{itemize}
\paragraph{Key Obligations}
\begin{itemize}[leftmargin=*,itemsep=0.2em]
\item Prove the source-to-object claim: the source-side content 'A barely-nontrivial lower bound for an explicit problem Q in a weak model' available at 'st.def.hardness-magnification-phenomenon-programme-leve' can be organized into 'MCSP witness family' as a family of instances equipped with a monotone witness measure.
\item Prove the payoff claim: once 'MCSP witness family' carries the route data from 'A barely-nontrivial lower bound for an explicit problem Q in a weak model' toward 'There exist general-purpose magnification theorems (not problem-specific) showing that certain mild lower bounds for...', it forces 'exponential-lower-bound' at 'st.thm.atserias-m-uller-simple-general-magnification-fr'.
\item Prove the transport claim: every admissible realization of 'MCSP witness family' preserves 'exponential-lower-bound witness complexity' while crossing the route nonuniform-circuits.
\item Exhibit the model-alignment claim: 'MCSP witness family' gives a faithful translation across nonuniform-circuits.
\end{itemize}
\paragraph{Main Risks}
\begin{itemize}[leftmargin=*,itemsep=0.2em]
\item Failure mode: 'MCSP witness family' never becomes more than topical overlap because the route-specific lemma remains missing. Prove that a hardness-magnification witness family upgrades the route from major separation consequence to exponential lower bound without losing 'exponential-lower-bound witness complexity'.
\item Abort logic: Abort if the negative-route intervention does not expose the expected failure: Allow 'MCSP witness family' to ignore the source-side content 'A barely-nontrivial lower bound for an explicit problem Q in a weak model' and only preserve topical overlap. => bridge-degrades-or-fails.
\item Global limitation: evidence\_count=0 and unresolved hypotheses=22, so the route is still under mathematical pressure.
\end{itemize}
\paragraph{Expert Actions}
\begin{itemize}[leftmargin=*,itemsep=0.2em]
\item Discharge the route-specific lemma: Prove that a hardness-magnification witness family upgrades the route from major separation consequence to exponential lower bound without losing 'exponential-lower-bound witness complexity'.
\item Stabilize the main bridge lemma: Source-to-object lemma: the source-side content 'A barely-nontrivial lower bound for an explicit problem Q in a weak model' at 'st.def.hardness-magnification-phenomenon-programme-leve' yields a realization of 'MCSP witness family' as a family of instances equipped with a monotone witness measure, satisfying admits source-side witness data for 'A barely-nontrivial lower bound for an explicit problem Q in a weak model'; preserves the route invariant 'exponential-lower-bound witness complexity'; aligns the model shift (nonuniform-circuits) with an explicit transport map.
\item Run the designated counterfactual probe: Allow 'MCSP witness family' to ignore the source-side content 'A barely-nontrivial lower bound for an explicit problem Q in a weak model' and only preserve topical overlap. bridge-degrades-or-fails The bridge object must absorb the actual source-side mathematical content rather than merely name the antecedent.
\end{itemize}

\section{Novelty Assessment}
\subsection{Positive Evidence}
\begin{itemize}[leftmargin=*,itemsep=0.2em]
\item The leading route is theorem-grounded at the level of source and target excerpts, not only labels: st.def.hardness-magnification-phenomenon-programme-leve $\rightarrow$  st.thm.oliveira-pich-santhanam-magnification-near-state.
\item The bridge objects are typed carriers such as 'MCSP witness family' with carrier 'family of instances equipped with a monotone witness measure', rather than generic placeholders.
\item The routes expose explicit missing lemmas and negative routes, so novelty is attached to a falsifiable mathematical program rather than to generic thematic similarity.
\end{itemize}

\subsection{Limits}
\begin{itemize}[leftmargin=*,itemsep=0.2em]
\item The current run still has evidence\_count=0, so support is structural and explanatory, not evidential.
\item There remain 22 inconclusive hypotheses, so the frontier is not yet mathematically settled.
\item The novelty is therefore programmatic: new route decompositions, typed carriers, and missing lemmas, not validated new theorems.
\end{itemize}

\subsection{Verdict}
There is real programmatic novelty in the sense of a new content-level bridge program: the framework isolates specific theorem-pair deltas, proposes typed intermediate carriers, and exposes missing route lemmas. There is not yet evidence of theorem-level novelty, because the decisive lemmas remain open and the run has not accumulated external or adversarial evidence.

\section{Next Best Direction To Explore}
\begin{itemize}[leftmargin=*,itemsep=0.25em]
\item Bridge ID: \texttt{bridge.002}
\item Title: MCSP witness family: major lower bounds for NP in much stronger models (including superpolynomial formula/circuit lower bounds, and in some parameter...
\item Route: \texttt{st.def.hardness-magnification-phenomenon-programme-leve} $\rightarrow$ \texttt{st.thm.oliveira-santhanam-hardness-magnification-for-na}
\item Assessment: This is the next best direction because it keeps the leading bridge object fixed while forcing it through a harder transport regime. If the same witness family survives this adjacent target, that is evidence that the carrier is mathematically real rather than tuned only to one theorem statement.
\item Missing route lemma: Prove that a hardness-magnification witness family preserves 'MCSP complexity profile in nonuniform-circuits' while transporting the route from the nonuniform-circuits regime to the formula-size regime strongly enough to derive 'major lower bounds for NP in much stronger models (including superpolynomial formula/circuit lower bounds, and in some parameter...'.
\end{itemize}

\section{Recommended Probe}
\begin{itemize}[leftmargin=*,itemsep=0.25em]
\item Hypothesis: \texttt{phyp.p-vs-np.021}
\item Target statement: \texttt{st.thm.oliveira-santhanam-hardness-magnification-for-na}
\item Priority: 0.690
\item Recommended action: run-another-abduce-deduce-round
\item Reason: high priority inherited from connector hypotheses
\end{itemize}

\end{document}
